% GENERATED FILE. Edit canonical lessons, then run npm run generate:books.
% canonical-source-hash: fnv1a64:fe42c6592673bae6
% canonical-lessons: TE-C382-alocana, TE-C382-drkpatham, TE-C382-abhiprayabhedam, TE-C382-bhavincu, TE-C382-konam

\chapter{Ideas and Points of View}
\label{ch:te-ideas-and-points-of-view}
\hlchaptermodality{382}

\begin{chapteropening}
\textbf{By the end of this chapter:} \emph{I can talk about ideas, outlooks and points of view.}

Complete the last lesson of chapter 382: I can talk about ideas, outlooks and points of view.
\end{chapteropening}

\section[\texorpdfstring{\te{ఆలోచన} (ālōcana)}{ఆలోచన (ālōcana)}]{\te{ఆలోచన} (ālōcana) — an idea; a thought}

\label{lesson:TE-C382-alocana}



\begin{quote}
Before the new one: say the Telugu for a watchman, then the Telugu for a milkman.
\end{quote}

\subsection*{You'll want to know: \te{ఆలోచన}}
\textbf{\te{ఆలోచన}} — \emph{ālōcana} — \textquotedblleft{}an idea; a thought\textquotedblright{}.

\begin{quote}
\te{నాకు} \te{ఒక} \te{మంచి} \te{ఆలోచన} \te{వచ్చింది}.
\end{quote}

\begin{quote}
\emph{nāku oka mañci ālōcana vaccindi.}
\end{quote}

\begin{quote}
I had a good idea.
\end{quote}

Plural: \te{ఆలోచనలు}, \emph{ālōcanalu}.

An idea comes to you: \te{నాకు} \te{ఆలోచన} \te{వచ్చింది}, literally an idea came to me.

\begin{cousinweb}
\textbf{\te{ఆలోచన}} (\emph{ālōcana}) is borrowed from Sanskrit \emph{ālocana}, \textquotedblleft{}seeing\textquotedblright{}, and also \textquotedblleft{}considering, reflection\textquotedblright{} (Monier-Williams).
\end{cousinweb}

\begin{grammarlens}[title={What you've built}]
One more word for giving reasons and opinions.
\end{grammarlens}

\subsection*{Your turn}
\begin{itemize}
\raggedright
  \item \emph{Say it:} \emph{ālōcana}
  \item \emph{Say it:} \emph{ālōcana}, once more
  \item \emph{Say it:} \emph{nāku oka mañci ālōcana vaccindi.} — I had a good idea.
  \item \emph{Recall:} say the Telugu for a watchman, then the Telugu for a milkman, then say \emph{ālōcana} again
\end{itemize}

\begin{tcolorbox}[breakable,colback=teal!4,colframe=teal!35!black,title={Before you move on}]
What does \te{ఆలోచన} mean? (\textquotedblleft{}An idea; a thought\textquotedblright{}.) Say it once more.
\end{tcolorbox}
\section[\texorpdfstring{\te{దృక్పథం} (dṛkpathaṁ)}{దృక్పథం (dṛkpathaṁ)}]{\te{దృక్పథం} (dṛkpathaṁ) — an outlook; a perspective}

\label{lesson:TE-C382-drkpatham}



\begin{quote}
Before the new one: say the Telugu for a milkman, then the Telugu for an idea.
\end{quote}

\subsection*{You'll want to know: \te{దృక్పథం}}
\textbf{\te{దృక్పథం}} — \emph{dṛkpathaṁ} — \textquotedblleft{}an outlook; a perspective\textquotedblright{}.

\begin{quote}
\te{ప్రతి} \te{మనిషికి} \te{ఒక} \te{దృక్పథం} \te{ఉంటుంది}.
\end{quote}

\begin{quote}
\emph{prati maniṣiki oka dṛkpathaṁ uṇṭundi.}
\end{quote}

\begin{quote}
Every person has an outlook.
\end{quote}

A formal word for the way someone looks at life or at a question.

\begin{cousinweb}
\textbf{\te{దృక్పథం}} (\emph{dṛkpathaṁ}) is borrowed from Sanskrit \emph{dṛkpatha}, \textquotedblleft{}the range of sight\textquotedblright{} (Monier-Williams), a compound of \emph{dṛk-}, from \emph{dṛś} \textquotedblleft{}to see\textquotedblright{}, and \emph{patha}, \textquotedblleft{}a path\textquotedblright{} (CDIAL 7743).
\end{cousinweb}

\begin{grammarlens}[title={What you've built}]
One more word for giving reasons and opinions.
\end{grammarlens}

\subsection*{Your turn}
\begin{itemize}
\raggedright
  \item \emph{Say it:} \emph{dṛkpathaṁ}
  \item \emph{Say it:} \emph{dṛkpathaṁ}, once more
  \item \emph{Say it:} \emph{prati maniṣiki oka dṛkpathaṁ uṇṭundi.} — Every person has an outlook.
  \item \emph{Recall:} say the Telugu for a milkman, then the Telugu for an idea, then say \emph{dṛkpathaṁ} again
\end{itemize}

\begin{tcolorbox}[breakable,colback=teal!4,colframe=teal!35!black,title={Before you move on}]
What does \te{దృక్పథం} mean? (\textquotedblleft{}An outlook; a perspective\textquotedblright{}.) Say it once more.
\end{tcolorbox}
\section[\texorpdfstring{\te{అభిప్రాయభేదం} (abhiprāyabhēdaṁ)}{అభిప్రాయభేదం (abhiprāyabhēdaṁ)}]{\te{అభిప్రాయభేదం} (abhiprāyabhēdaṁ) — difference of opinion}

\label{lesson:TE-C382-abhiprayabhedam}



\begin{quote}
Before the new one: say the Telugu for an idea, then the Telugu for an outlook.
\end{quote}

\subsection*{You'll want to know: \te{అభిప్రాయభేదం}}
\textbf{\te{అభిప్రాయభేదం}} — \emph{abhiprāyabhēdaṁ} — \textquotedblleft{}difference of opinion\textquotedblright{}.

\begin{quote}
\te{మా} \te{ఇద్దరి} \te{మధ్య} \te{చిన్న} \te{అభిప్రాయభేదం} \te{వచ్చింది}.
\end{quote}

\begin{quote}
\emph{mā iddari madhya cinna abhiprāyabhēdaṁ vaccindi.}
\end{quote}

\begin{quote}
A small difference of opinion came up between the two of us.
\end{quote}

Two words you can take apart: \te{అభిప్రాయం}, an opinion, which you know, and \te{భేదం}, a difference.

\begin{cousinweb}
\textbf{\te{అభిప్రాయభేదం}} (\emph{abhiprāyabhēdaṁ}) is built from two Sanskrit words: \emph{abhiprāya}, \textquotedblleft{}aim, intention\textquotedblright{}, which Telugu uses for an opinion, and \emph{bheda}, \textquotedblleft{}a splitting, a breaking apart\textquotedblright{} (Monier-Williams).
\end{cousinweb}

\begin{grammarlens}[title={What you've built}]
One more word for giving reasons and opinions.
\end{grammarlens}

\subsection*{Your turn}
\begin{itemize}
\raggedright
  \item \emph{Say it:} \emph{abhiprāyabhēdaṁ}
  \item \emph{Say it:} \emph{abhiprāyabhēdaṁ}, once more
  \item \emph{Say it:} \emph{mā iddari madhya cinna abhiprāyabhēdaṁ vaccindi.} — A small difference of opinion came up between the two of us.
  \item \emph{Recall:} say the Telugu for an idea, then the Telugu for an outlook, then say \emph{abhiprāyabhēdaṁ} again
\end{itemize}

\begin{tcolorbox}[breakable,colback=teal!4,colframe=teal!35!black,title={Before you move on}]
What does \te{అభిప్రాయభేదం} mean? (\textquotedblleft{}Difference of opinion\textquotedblright{}.) Say it once more.
\end{tcolorbox}
\section[\texorpdfstring{\te{భావించు} (bhāviñcu)}{భావించు (bhāviñcu)}]{\te{భావించు} (bhāviñcu) — to regard as; to consider}

\label{lesson:TE-C382-bhavincu}



\begin{quote}
Before the new one: say the Telugu for an outlook, then the Telugu for difference of opinion.
\end{quote}

\subsection*{You'll want to know: \te{భావించు}}
\textbf{\te{భావించు}} — \emph{bhāviñcu} — \textquotedblleft{}to regard as; to consider\textquotedblright{}.

\begin{quote}
\te{నేను} \te{ఆ} \te{పనిని} \te{నా} \te{బాధ్యతగా} \te{భావించాను}.
\end{quote}

\begin{quote}
\emph{nēnu ā panini nā bādhyatagā bhāviñcānu.}
\end{quote}

\begin{quote}
I regarded that work as my responsibility.
\end{quote}

Past (I …): \te{భావించాను}, \emph{bhāviñcānu}.

The thing takes -\te{ని} and what you take it to be takes -\te{గా}: \te{దాన్ని} \te{బాధ్యతగా} \te{భావించు}.

\begin{cousinweb}
\textbf{\te{భావించు}} (\emph{bhāviñcu}) is built on Sanskrit \emph{bhāvana}, \textquotedblleft{}causing to be\textquotedblright{}, and as a noun \textquotedblleft{}imagination\textquotedblright{} (CDIAL 9476), from the root \emph{bhū}, \textquotedblleft{}to be\textquotedblright{}, with the Telugu verb ending -\te{ఇంచు}.
\end{cousinweb}

\begin{grammarlens}[title={What you've built}]
One more word for giving reasons and opinions.
\end{grammarlens}

\subsection*{Your turn}
\begin{itemize}
\raggedright
  \item \emph{Say it:} \emph{bhāviñcu}
  \item \emph{Say it:} \emph{bhāviñcu}, once more
  \item \emph{Say it:} \emph{nēnu ā panini nā bādhyatagā bhāviñcānu.} — I regarded that work as my responsibility.
  \item \emph{Recall:} say the Telugu for an outlook, then the Telugu for difference of opinion, then say \emph{bhāviñcu} again
\end{itemize}

\begin{tcolorbox}[breakable,colback=teal!4,colframe=teal!35!black,title={Before you move on}]
What does \te{భావించు} mean? (\textquotedblleft{}To regard as; to consider\textquotedblright{}.) Say it once more.
\end{tcolorbox}
\section[\texorpdfstring{\te{కోణం} (kōṇaṁ)}{కోణం (kōṇaṁ)}]{\te{కోణం} (kōṇaṁ) — an angle; a point of view}

\label{lesson:TE-C382-konam}



\begin{quote}
Before the new one: say the Telugu for difference of opinion, then the Telugu for to regard as.
\end{quote}

\subsection*{You'll want to know: \te{కోణం}}
\textbf{\te{కోణం}} — \emph{kōṇaṁ} — \textquotedblleft{}an angle; a point of view\textquotedblright{}.

\begin{quote}
\te{మనం} \te{ఆ} \te{సమస్యను} \te{వేరే} \te{కోణం} \te{నుంచి} \te{చూడాలి}.
\end{quote}

\begin{quote}
\emph{manaṁ ā samasyanu vērē kōṇaṁ nuñci cūḍāli.}
\end{quote}

\begin{quote}
We should view that problem from a different angle.
\end{quote}

Plural: \te{కోణాలు}, \emph{kōṇālu}.

In geometry a corner or an angle; in talk, the side from which you see a question.

\begin{cousinweb}
\textbf{\te{కోణం}} (\emph{kōṇaṁ}) is borrowed from Sanskrit \emph{koṇa}, \textquotedblleft{}a corner, an angle\textquotedblright{} (Monier-Williams). Like many Sanskrit nouns in -a, it ends in -\te{ం} in Telugu.
\end{cousinweb}

\begin{grammarlens}[title={What you've built}]
One more word for giving reasons and opinions.
\end{grammarlens}

\subsection*{Your turn}
\begin{itemize}
\raggedright
  \item \emph{Say it:} \emph{kōṇaṁ}
  \item \emph{Say it:} \emph{kōṇaṁ}, once more
  \item \emph{Say it:} \emph{manaṁ ā samasyanu vērē kōṇaṁ nuñci cūḍāli.} — We should view that problem from a different angle.
  \item \emph{Recall:} say the Telugu for difference of opinion, then the Telugu for to regard as, then say \emph{kōṇaṁ} again
\end{itemize}

\begin{tcolorbox}[breakable,colback=teal!4,colframe=teal!35!black,title={Before you move on}]
What does \te{కోణం} mean? (\textquotedblleft{}An angle; a point of view\textquotedblright{}.) Say it once more.
\end{tcolorbox}
